\documentclass[aspectratio=169,11pt]{beamer}
% Shared Beamer style for practice contest solution walkthroughs.
\usepackage[utf8]{inputenc}
\usepackage[T1]{fontenc}
\usepackage{lmodern}
\usepackage{amsmath,amssymb}
\usepackage{xcolor}
\usepackage{hyperref}
\usepackage{listings}

\definecolor{CPNavy}{HTML}{172033}
\definecolor{CPBlue}{HTML}{2563EB}
\definecolor{CPGreen}{HTML}{16A34A}
\definecolor{CPOrange}{HTML}{F97316}
\definecolor{CPGray}{HTML}{64748B}
\definecolor{CPLight}{HTML}{F8FAFC}
\definecolor{CPCodeBg}{HTML}{F1F5F9}
\definecolor{CPCodeRule}{HTML}{CBD5E1}

\usetheme{default}
\usefonttheme{professionalfonts}
\setbeamertemplate{navigation symbols}{}
\setbeamersize{text margin left=0.65cm,text margin right=0.65cm}
\setbeamertemplate{itemize item}{\color{CPBlue}$\blacktriangleright$}
\setbeamertemplate{itemize subitem}{\color{CPGray}$\triangleright$}

\setbeamercolor{normal text}{fg=CPNavy,bg=white}
\setbeamercolor{frametitle}{fg=white,bg=CPNavy}
\setbeamercolor{title}{fg=white,bg=CPNavy}
\setbeamercolor{subtitle}{fg=CPLight,bg=CPNavy}
\hypersetup{colorlinks=true,urlcolor=CPBlue}

\setbeamertemplate{frametitle}{%
  \nointerlineskip%
  \begin{beamercolorbox}[wd=\paperwidth,ht=0.88cm,dp=0.22cm,leftskip=0.55cm,rightskip=0.35cm]{frametitle}
    \usebeamerfont{frametitle}\insertframetitle\hfill\small\insertframenumber/\inserttotalframenumber
  \end{beamercolorbox}%
}

\setbeamertemplate{footline}{%
  \leavevmode%
  \hbox{%
    \begin{beamercolorbox}[wd=.58\paperwidth,ht=2.4ex,dp=1ex,leftskip=.5cm]{author in head/foot}%
      \color{CPGray}\insertshorttitle
    \end{beamercolorbox}%
    \begin{beamercolorbox}[wd=.42\paperwidth,ht=2.4ex,dp=1ex,rightskip=.5cm plus1fil]{date in head/foot}%
      \hfill\color{CPGray}\insertshortauthor
    \end{beamercolorbox}%
  }%
}

\setbeamertemplate{title page}{%
  \vspace*{1.25cm}
  \begin{beamercolorbox}[wd=\paperwidth,sep=0.65cm,leftskip=0.15cm,rightskip=0.15cm]{title}
    {\color{white}\Huge\bfseries\inserttitle\par}
    \vspace{0.35cm}
    {\color{CPLight}\Large\insertsubtitle\par}
    \vspace{0.6cm}
    {\color{CPLight}\large\insertauthor\par}
  \end{beamercolorbox}
  \vfill
}

\newcommand{\sectiontag}[1]{\textbf{\color{CPBlue}#1}}
\newenvironment{tightitemize}{\begin{itemize}\small\setlength{\itemsep}{0.18cm}}{\end{itemize}}
\lstdefinestyle{cpcode}{
  language=Python,
  basicstyle=\ttfamily\tiny,
  keywordstyle=\color{CPBlue}\bfseries,
  commentstyle=\color{CPGray}\itshape,
  stringstyle=\color{CPGreen},
  backgroundcolor=\color{CPCodeBg},
  frame=single,
  framerule=0.25pt,
  rulecolor=\color{CPCodeRule},
  showstringspaces=false,
  columns=fullflexible,
  keepspaces=true,
  breaklines=true,
  breakatwhitespace=false,
  tabsize=4,
  xleftmargin=0.08cm,
  xrightmargin=0.08cm,
  aboveskip=0.08cm,
  belowskip=0.08cm
}


\title[enlarginghashtables]{Enlarging Hash Tables}
\subtitle{Day 5 solution walkthrough}
\author[Solution Deck]{Competitive Programming Bootcamp}
\date{}

\begin{document}

\begin{frame}[plain]
  \titlepage
\end{frame}

% BEGIN GENERATED STATEMENT INTERPRETATION
\begin{frame}{Interpret the Word Problem}
\small
\sectiontag{Statement excerpts:}
\begin{tightitemize}
  \item \textit{``smallest prime number larger than \$2n\$''}
  \item \textit{``if \$n\$ happens not to be prime''}
\end{tightitemize}

\vspace{0.18cm}
\sectiontag{Programming interpretation:}
\begin{tightitemize}
  \item For each input n, search upward from 2n + 1 until a prime is found.
  \item Also test whether the old table size n is prime.
  \item The output changes only by adding a parenthetical warning for composite n.
\end{tightitemize}
\end{frame}
% END GENERATED STATEMENT INTERPRETATION







\begin{frame}{Problem Model}
\small
\sectiontag{Textbook connection:} CP4 5.3.1 Prime Numbers

\vspace{0.25cm}
\sectiontag{Core technique:} Primality testing by trial division

\vspace{0.25cm}
\begin{tightitemize}
  \item For each n, find the smallest prime strictly greater than 2n.
  \item Also report if the original n is not prime.
  \item Input ends at n = 0.
\end{tightitemize}
\end{frame}

\begin{frame}{How to Recognize the Approach}
\begin{tightitemize}
  \item A candidate larger than 2n can be tested independently for primality.
  \item After 2, only odd candidates need to be checked.
  \item Trial division up to sqrt(n) is enough for the input scale.
\end{tightitemize}
\end{frame}

\begin{frame}{Algorithm}
\begin{tightitemize}
  \item Implement is\_prime(n).
  \item Reject n < 2, accept n = 2, reject other even numbers.
  \item Try odd divisors from 3 through isqrt(n).
  \item For each input n, start candidate = 2n + 1 and advance by 2 until prime.
  \item Print candidate, adding the warning if n itself is not prime.
\end{tightitemize}
\end{frame}

\begin{frame}{Walkthrough of the Key State}
\begin{tightitemize}
  \item If n has a divisor above sqrt(n), the paired divisor is below sqrt(n).
  \item 2n + 1 is always odd, so candidate += 2 preserves odd candidates.
  \item The original n and enlarged candidate are separate primality checks.
\end{tightitemize}
\end{frame}

\begin{frame}{Why the Algorithm Is Correct}
\begin{tightitemize}
  \item is\_prime tests all possible factor pairs by checking divisors up to the square root.
  \item The candidate loop checks increasing integers greater than 2n, skipping only even composites.
  \item The first prime found is therefore the smallest valid enlarged table size.
\end{tightitemize}
\end{frame}

\begin{frame}{Complexity and Implementation Notes}
\small
\sectiontag{Complexity:} O(sqrt(C)) per primality test, where C is the candidate value; O(1) memory.

\vspace{0.3cm}
\begin{tightitemize}
  \item Use math.isqrt to avoid floating-point square-root issues.
  \item Remember the exact output format for non-prime original n.
\end{tightitemize}
\end{frame}



% BEGIN GENERATED CODE WALKTHROUGH
\begin{frame}[fragile]{Code: Primality Test}
\scriptsize
\sectiontag{Source lines:} \texttt{enlarginghashtables.py:43--70}

\vspace{0.08cm}
\begin{columns}[T,totalwidth=\textwidth]
\begin{column}{0.58\textwidth}
\begin{lstlisting}[style=cpcode]
import sys
from math import isqrt

def is_prime(n):
    if n < 2:
        return False

    if n == 2:
        return True

    if n % 2 == 0:
        return False

    limit = isqrt(n)

    for divisor in range(3, limit + 1, 2):
        if n % divisor == 0:
            return False

    return True
\end{lstlisting}
\end{column}
\begin{column}{0.38\textwidth}
\sectiontag{What this chunk does}
\begin{tightitemize}
  \item Numbers below 2 and even numbers above 2 are handled immediately.
  \item Only odd divisors up to sqrt(n) need to be tested.
  \item If no divisor is found, n is prime.
\end{tightitemize}
\end{column}
\end{columns}
\end{frame}

\begin{frame}[fragile]{Code: Choose First Candidate}
\scriptsize
\sectiontag{Source lines:} \texttt{enlarginghashtables.py:73--85}

\vspace{0.08cm}
\begin{columns}[T,totalwidth=\textwidth]
\begin{column}{0.58\textwidth}
\begin{lstlisting}[style=cpcode]
def main():
    input = sys.stdin.buffer.readline

    while True:
        n = int(input())

        if n == 0:
            break

        candidate = 2 * n + 1
\end{lstlisting}
\end{column}
\begin{column}{0.38\textwidth}
\sectiontag{What this chunk does}
\begin{tightitemize}
  \item Input is processed until the sentinel value 0.
  \item The first possible new size is strictly greater than twice n.
  \item 2n + 1 is odd, so later candidates can advance by 2.
\end{tightitemize}
\end{column}
\end{columns}
\end{frame}

\begin{frame}[fragile]{Code: Find and Print New Size}
\scriptsize
\sectiontag{Source lines:} \texttt{enlarginghashtables.py:87--101}

\vspace{0.08cm}
\begin{columns}[T,totalwidth=\textwidth]
\begin{column}{0.58\textwidth}
\begin{lstlisting}[style=cpcode]
        while not is_prime(candidate):
            candidate += 2

        if is_prime(n):
            print(candidate)

        else:
            print(f"{candidate} ({n} is not prime)")

main()
\end{lstlisting}
\end{column}
\begin{column}{0.38\textwidth}
\sectiontag{What this chunk does}
\begin{tightitemize}
  \item The loop stops at the first prime candidate, making it the smallest valid new size.
  \item The original n is tested separately.
  \item Composite original sizes receive the required note.
\end{tightitemize}
\end{column}
\end{columns}
\end{frame}
% END GENERATED CODE WALKTHROUGH

\begin{frame}{Source}
\small
\sectiontag{Solution file:} \texttt{practice\_contests/day\_5/enlarginghashtables.py}

\vspace{0.3cm}
\sectiontag{Problem:} \url{https://open.kattis.com/problems/enlarginghashtables}

\vspace{0.45cm}
Use this deck with the source code open beside it. The slides explain the reasoning; the code shows the exact input handling and output formatting.
\end{frame}

\end{document}
